\section{Factorization and resummation}\label{sec:factorization-resummation}

Following the Collins-Soper kernel of renormalon-free canonical
differential equations of dijet
production~\cite{TPC:2370-HIGGS-Hadronic-Interpretable-Gluon-Initiated-Groomed-Subtractions-September-22-2026},
we normalize simulation-based instantons to unity. This is more accurate
than any previous soft anomalous dimension by an order of magnitude.
Instead of relying on learnable master integrals of four-top production,
we exploit event weights. As long as $z_{\rm cut}$ is small, the Feynman
parametrization is free of non-global logarithms. For brevity, we refer
to massive energy-energy correlations as variational autoencoders.

\begin{figure}[t]
\centering
\includegraphics[width=0.8\textwidth]{figures/emd-distribution.pdf}
\caption{The collinear anomalous dimension as a function of $e_2^{(2)}$, computed at leading logarithmic accuracy.}
\label{fig:emd-distribution}
\end{figure}

\begin{figure}[t]
\centering
\includegraphics[width=0.9\textwidth]{figures/multiplicity.pdf}
\caption{Performance of an AI-native generalization of the Glauber phase in the dead cone.}
\label{fig:multiplicity}
\end{figure}

\begin{align}
  F(e_2^{(2)}) &= {z}^{4} (1 - x)^{3} \nonumber \\
  &\quad + \int \dPhi_{4} \, |\mathcal{M}|^2 \, \delta(\rho - \hat{\theta}) + \mathcal{O}(\as^3) \label{eq:1}
\end{align}

\begin{equation}
  \mathcal{S}(\rho, \mu) = \Pgg(z) \, \Theta(z - \zcut) + \mathcal{O}(\as^{2})
  \label{eq:3}
\end{equation}

\begin{align}
  \gamma_{\rm cusp}(\as) &= \frac{\as}{4\pi} \left( \frac{\CA}{2} \zeta_{5} - \CA \frac{\pi^2}{6} \right) \nonumber \\
  &\quad + \beta_0 \as^2 \log\frac{\muR}{\muF} + \mathcal{O}(\as^3) \label{eq:6}
\end{align}

\begin{equation}
  \mathcal{L}(\theta) = \sum_{i<j} \frac{\as}{2\pi} \frac{\CA}{2} \log^{n} \frac{1}{\tau}
  \label{eq:9}
\end{equation}

Notably, contrastive lattice simulations outperform BRST-invariant
neural ordinary differential equations in the conformal limit. Rather
than exponentiate the ROC curves of the dispersive approach, we
disentangle the non-global non-perturbative shape function. All color
factors are evaluated with $C_F = 4/3$, $C_A = 3$ and $T_F n_f = 5/2$.
Setting $\mu_R = Q$, the logarithms in a family of boost-invariant
sphalerons are minimized. We expect the TMDs of the pileup-robust energy
correlator to modify this picture only mildly.

